\begin{longtable}[c]{|l|m{10cm}|}
    \caption{移位指令}
    \label{tab:pie-shift-instructions}\\
    \hline
    \rowcolor{lightgray}
    \textbf{指令} & \textbf{定义} \\\hline
    \endfirsthead
    \rowcolor{lightgray}
    \textbf{指令} & \textbf{定义} \\\hline
    \endhead
    \endfoot

    \hyperref[instruction:ESPSRCQ]{ESP.SRC.Q} & 16-字节拼接后逻辑右移，移位数值由 SAR\_BYTE 寄存器决定。
    \\\hline
\hyperref[instruction:ESPSRCQLDXP]{ESP.SRC.Q.LD.XP} & 对拼接的16字节数据执行逻辑右移，移位值由SAR\_BYTE寄存器决定。与此同时，从内存读取16字节数据，并将寄存器值加到读取地址上。 \\\hline
\hyperref[instruction:ESPSRCQLDIP]{ESP.SRC.Q.LD.IP} & 对拼接的16字节数据执行逻辑右移，移位值由SAR\_BYTE寄存器决定。与此同时，从内存读取16字节数据，并将立即数加到读取地址上。 \\\hline
\hyperref[instruction:ESPSLCI2Q]{ESP.SLCI.2Q} & 对拼接的16字节数据执行逻辑左移，移位值由立即数决定。 \\\hline
\hyperref[instruction:ESPSLCXXP2Q]{ESP.SLCXXP.2Q} & 对拼接的16字节数据执行逻辑左移，移位值由AR寄存器中的值决定。 \\\hline
\hyperref[instruction:ESPSRCI2Q]{ESP.SRCI.2Q} & 对拼接的16字节数据执行逻辑右移，移位值由立即数决定。 \\\hline
\hyperref[instruction:ESPSRCXXP2Q]{ESP.SRCXXP.2Q} & 对拼接的16字节数据执行逻辑右移，移位值由AR寄存器中的值决定。 \\\hline
\hyperref[instruction:ESPSRCQ128STINCP]{ESP.SRCQ.128.ST.INCP} & 对拼接的16字节数据执行逻辑右移，移位后的数据将被写入内存。 \\\hline
\hyperref[instruction:ESPVSRU32]{ESP.VSR.U32} & 对4字节数据执行向量算术右移。 \\\hline
\hyperref[instruction:ESPVSRS32]{ESP.VSR.S32} & 对4字节数据执行向量算术右移。 \\\hline
\hyperref[instruction:ESPVSL32]{ESP.VSL.32} & 对4字节数据执行向量算术左移。 \\\hline
\hyperref[instruction:ESPVSrd32]{ESP.VSRD.[8/16/32]} & 对1字节/2字节/4字节数据执行向量算术右移或左移。 \\\hline
\hyperref[instruction:ESPVSLD32]{ESP.VSLD.[8/16/32]} & 对1字节/2字节/4字节数据执行向量算术右移或左移。 \\\hline

\end{longtable}